% Bản nháp
\documentclass[../main.tex]{subfiles}
\begin{document}
%==========================================TIÊU ĐỀ TẬP========================================
\begin{table}[h]
    \begin{adjustwidth}{-1cm}{}
        \begin{tabular}{>{\centering\arraybackslash}p{6cm}|>{\raggedright\arraybackslash}p{12.5cm}}
            \multirow{5}{6cm}
            % Cột bên trái: ÉPISODE và số tập
            {\\[-40pt]
                \flaregothic\fontsize{20pt}{20pt}\selectfont \centering
                \color{eptype}SPÉCIAL\color{black}\\
                \burbank\fontsize{72pt}{72pt}\selectfont
                \color{epnum}XVII\color{black}
                \\[18pt]
                \flaregothic\fontsize{18pt}{18pt}\selectfont
                \color{epattr}KỲ THI VÀO 10\\
                \color{black}
            }
             &
            % Dòng thứ nhất: tên tập tiếng Nhật
            {\fotskip\fontsize{18pt}{18pt}\selectfont \ruby{安}{あん}\ruby{和}{わ}\ruby{高}{こう}\ruby{校}{こう}への\ruby{道}{みち}
            }                                    \\[6pt]
            % Dòng thứ hai: tên tập tiếng Anh
             & {\cafeteria\fontsize{18pt}{18pt}\selectfont The Gateway to Anwa}                                                                       \\[4pt]
            % Dòng thứ ba: tên tập tiếng Việt
             & {\vnmsans\fontsize{18pt}{18pt}\selectfont Cánh cổng Yên Hòa}                                                                           \\[8pt]
            % Dòng thứ tư: Nhân vật xuất hiện
             & {\caviar{\fontsize{14pt}{14pt}\selectfont Nhân vật: \emp{kuon2} \emp{kaigou2} \emp{suzuki2} \emp{ryuuhou2} \emp{naomi2} \emp{mizuha}}}
            \\[8pt]
            % Dòng cuối cùng: Thời gian diễn ra của tập (thường sẽ là ngày phát hành)
            %    & {\continuum\fontsize{16pt}{16pt}\selectfont Thời gian: Tuyển sinh năm học mới}                                                    \\[18pt]
        \end{tabular}
    \end{adjustwidth}
\end{table}
%=========================================TRUYỆN CHÍNH========================================
\color{black}

\txtbx[2]{{\color{vol} \uline{Tóm tắt:}}} Bước vào giai đoạn tuyển sinh lớp 10 công lập khốc liệt tại Kawachi, ba cô em gái nhà Hikawa gồm Suzuki, Ryuuhou và Miyako (dưới danh nghĩa Mizuha) đồng loạt đặt nguyện vọng 1 vào ngôi trường điểm danh tiếng Anwa (Yên Hòa). Đối diện với sức nóng nghẹt thở của kỳ thi, cả gia đình 6 anh chị em cùng cô bé Naomi đã phải chung tay lập chiến dịch cứu trợ học tập cấp tốc. Chương truyện mở ra những góc nhìn hài hước về chặng đường ôn thi đầy nước mắt của Ryuuhou và giai thoại đỗ thủ khoa đầy ngẫu nhiên của người anh cả Kuon.

Mùa hè năm ấy, không khí tại nội thành Kawachi trở nên ngột ngạt và nóng bức hơn bao giờ hết. Nhưng cái nóng của thời tiết hoàn toàn chẳng thấm tháp vào đâu so với bầu không khí căng thẳng đang đè nặng lên ngôi nhà của gia đình Hikawa.

Trên chiếc bàn gỗ lớn ở phòng khách, sách giáo khoa, các tập đề cương ôn tập môn Toán, Ngữ văn, Tiếng Anh và Tiếng Nhật chất cao như những bức tường thành nhỏ. Ba cô em gái — Suzuki, Ryuuhou và Miyako — chuẩn bị bước vào kỳ thi tuyển sinh lớp 10 công lập, một trong những kỳ thi được đánh giá là khốc liệt và cam go bậc nhất tại Etsunan. Và mục tiêu chung mà cả ba nhắm tới chính là Anwa, ngôi trường điểm công lập danh tiếng hàng đầu quận Kyuukami.

``Aaa! Cái đống hình học không gian này bị cái lỗi hệ thống gì thế không biết! Sao cái đường thẳng này nhìn kiểu gì cũng không cắt được mặt phẳng thế!?'' Tiếng gào thét đầy cọc cằn quen thuộc của Ryuuhou vang lên xé toạc sự tĩnh lặng. Con bé vò đầu bứt tai, hai mắt đỏ ngầu nhìn chằm chằm vào hình vẽ hình chóp trên trang vở nháp. Khi áp lực học tập tăng cao, cái nhân cách ``game thủ toxic'' của con bé bỗng chốc được kích hoạt ngay trong lúc làm bài tập về nhà.

Kuon đang ngồi bên cạnh kiểm tra lại đống tài liệu, khẽ thở dài thườn thượt rồi vỗ vỗ vai em gái dỗ dành: ``Bớt nóng, bớt nóng đi em gái. Em vẽ lệch mất cái giao tuyến thế kia thì có đến mùa quýt năm sau nó mới chịu cắt nhau được nhé. Tập trung hít thở sâu, ăn viên kẹo ngọt này vào cho hạ hỏa đi em.''

``Học hành kiểu này thì trượt từ vòng gửi xe chứ đòi mơ tưởng đặt nguyện vọng vào trường điểm Yên Hòa hả cưng?'' Kaigou đứng tựa lưng vào thành cửa, khoanh hai tay trước ngực, khẽ buông lời nhận xét vô cùng sắc sảo và phũ phàng. ``Aniki hồi trước dẫu có lơ là thật nhưng ít ra cái đầu óc đại số của anh ấy còn có chút logic, chứ nhìn cái cách giải toán cơ bắp này của mày thì chị e là không ổn chút nào đâu.''

Ryuuhou nghe thấy thế liền phồng hai má lên giận dỗi quay phắt sang lườm chị gái: ``Chị nói thế mà nghe được à! Em hùa theo hai người đặt nguyện vọng vào Anwa là vì muốn được học tập dưới mái trường ngày xưa của ông anh Kuon thôi! Làm gì mà dìm hàng dập tắt hy vọng của em kinh thế chứ!''

Trái ngược hoàn toàn với sự náo loạn của Ryuuhou, Suzuki lại đang ngồi vô cùng ngăn nắp, cặm cụi nêm nếm chép lại từng công thức ngữ pháp Tiếng Nhật vào cuốn sổ tay nhỏ. Suzuki có lực học khá giỏi, tính cách lại vô cùng chăm chỉ, cần cù chịu khó để bù đắp lại những khoảng trống kiến thức. Cô bé ngẩng mặt lên cười hiền dịu: ``Onee-sama đừng trêu Ryuuhou nữa. Em tin là nếu tụi mình cùng nhau cố gắng ôn luyện nghiêm túc thì kiểu gì cũng sẽ có cơ hội đỗ đạt thôi mà.''

Mizuha (Miyako) ngồi ở góc bàn bên cạnh, điệu bộ khép népch thiệp đúng chuẩn phong thái của một thiếu nữ loài người, mái tóc đen mượt che giấu hoàn toàn đôi tai thú bấy giờ khẽ lay động nhẹ. Trước mặt cô nàng là một xấp đề thi thử nâng cao môn Toán được giải quyết sạch sẽ gọn gàng, không sót một câu nào với số điểm tối đa. Với trí tuệ siêu việt được buff mạnh mẽ gấp 10 lần nhờ tác dụng phụ từ thí nghiệm của Koyori, cái kỳ thi lớp 10 này đối với cô nàng thực sự chỉ là một trò chơi trẻ con vô vị. Mizuha khẽ đưa mắt cảnh giác nhìn Kaigou, rồi dịu dàng quay sang Suzuki: ``Suzuki-sama cứ yên tâm học tập đi ạ. Em nhất định sẽ luôn kề vai sát cánh hỗ trợ và bảo vệ chị vượt qua kỳ thi này.''

Thực chất, lý do Mizuha quyết tâm thi vào Anwa vô cùng đơn giản: cô nàng lo sợ sự nghiêm khắc bảo bọc có phần cực đoan của cô chị lớn Kaigou sẽ làm tổn hại đến Suzuki, nên muốn đích thân đi theo làm lá chắn bảo vệ ngầm cho cô chủ nhỏ của mình.

Cô bé Naomi 10 tuổi từ trong bếp thong thả bưng ra một khay nước trà đào mát lạnh cùng đĩa bánh bao ngọt thơm phức, cất giọng nói trong trẻo như tiếng chuông khánh: ``Các anh chị nghỉ tay uống chút nước giải khát cho mát mẻ tâm trí đi ạ. Naomi có chuẩn bị bánh bao ngọt để tiếp thêm năng lượng cho mọi người đây.''

Sự hiện diện bình yên và vầng hào quang ấm áp, tinh khiết của cô bé Naomi lập tức lan tỏa khắp căn phòng, xoa dịu hoàn toàn đi cái sự nóng bức căng thẳng của buổi học chiều. Ryuuhou vừa nhận lấy cốc trà đào từ tay Naomi liền hít một hơi thật sâu, gò má khẽ đỏ bừng ngượng ngùng, rũ bỏ hoàn toàn cái nhân cách toxic ban nãy, ngoan ngoãn cúi đầu: ``Oa, chị cảm ơn Naomi-chan nhiều nhiều lắm nha\~{} ngọt ngào quá đi mất hà!''

Kuon khẽ nhấp một ngụm trà, ngả người tựa lưng ra thành ghế, đưa mắt nhìn đống tài liệu rồi khẽ thở dài cười khổ: ``Thú thật là nhìn cái khoảng cách giữa năng lực học tập thực tế của Ryuuhou và cái điểm chuẩn đầu vào cao ngất ngưởng của trường mà em chọn lúc này, anh cảm thấy có một sự quan ngại sâu sắc không hề nhẹ đâu nhé. Suzuki thì còn có cửa nỗ lực vớt vát, Mizuha thì anh chắc chắn trăm phần trăm là đỗ thủ khoa rồi không cần bàn cãi, còn riêng em út Ryuuhou thì anh lo sốt vó luôn đây này.''

Ryuuhou bĩu môi hậm hực: ``Ông anh cứ lo hão! Ngày xưa ông anh cũng lười chảy thây ra mà vẫn đỗ đạt Yên Hòa đấy thôi, kể lại giai thoại cho tụi em nghe lấy động lực xem nào!''

Kuon nghe vậy liền bật cười chát chúa, nhớ lại ký ức thời trẻ trâu của mình: ``Haha, ngày xưa anh thi vào Anwa á? Nói thật là cái kịch bản lúc đó nó đi vào lòng đất lắm em ơi. Hồi đó anh hoàn toàn chẳng hề có một chút khái niệm hay nhận thức gì về cái sức nóng khốc liệt của kỳ thi cấp 3 này cả, cứ đinh ninh nghĩ nó cũng bình thường như mấy bài kiểm tra học kỳ ở trường cấp 2 thôi. Trong khi bạn bè cùng trang lứa học ngày học đêm, cắm đầu cắm cổ vào các lò luyện thi thêm bên ngoài sấp mặt, thì anh đây vẫn cứ thong thả bình chân như vại, ngày ngày sinh hoạt ung dung tự tại, đi học về là ôm máy tính chơi game như thể cái kỳ thi tuyển sinh kia không có liên quan gì đến cuộc đời mình vậy.''

Kaigou khẽ tặc lưỡi bóc mẽ: ``Chuẩn luôn đấy, cái đợt đó Aniki làm cho cả nhà được một phen lo sốt vó, đứng ngồi không yên luôn cơ.''

Kuon cười trừ tiếp tục kể: ``Thì tại vì anh đâu có biết Anwa thực chất lại là một trong những ngôi trường điểm danh tiếng top đầu của thành phố Kawachi đâu chứ! Cái lý do duy nhất khiến anh nhắm mắt đưa chân đặt nguyện vọng 1 vào trường đó thực ra rất là\dots\ ba chấm: đơn giản là vì cái trường đó nó cách nhà mình rất gần, đi bộ có mươi phút là tới nơi rồi, không phải tốn công thức dậy sớm chạy xe buýt chen chúc làm gì cho mệt người. Phải mãi cho đến tận cái ngày nhà trường công bố bảng điểm tổng kết và điểm chuẩn đầu vào, nhìn thấy cái tên mình nằm chễm chệ ở danh sách đỗ sát nút, anh mới hoàn toàn ngã ngửa kinh hoàng sực nhận ra bản thân vừa mới trải qua một cái kỳ thi cam go và khốc liệt nhất trong cuộc đời, thậm chí xét về độ cạnh tranh tỉ lệ chọi còn khó nuốt hơn cả kỳ thi đại học sau này nhiều lần cơ đấy.''

Nghe xong giai thoại thi cử đầy rẫy sự ngẫu nhiên của ông anh cả, cả Suzuki lẫn Ryuuhou đều phải tròn mắt kinh ngạc, còn Mizuha thì vẫn giữ ánh mắt lặng lẽ quan sát đầy trung thành.

Kuon đanh mặt lại, nghiêm túc gõ ngón tay xuống bàn họp: ``Thế nên anh mới bảo là ăn may thì chỉ có một lần thôi em ạ. Chứ còn đối với cái tình hình học lực bết bát hiện tại của em thì không có chuyện thần may mắn mỉm cười lần thứ hai đâu nhé, Ryuuhou. Từ ngày hôm nay trở đi, toàn thể 5 anh chị em trong nhà chúng ta sẽ cùng nhau lập chiến dịch 'chăm sóc' đặc biệt thiết quân luật uốn nắn học tập cho em, quyết tâm đập tan cái bức tường thi cử này mới thôi!''

``Trời đất ơi cứu tôi với Naomi-chan ơi! Ông anh già bắt đầu dở chứng bạo lực gia đình hành hạ em gái rồi đây này!'' Ryuuhou mếu máo ôm lấy cánh tay Naomi cầu cứu đầy kịch tính, khiến cho cả căn phòng lại được một phen cười rộ lên vui vẻ ấm áp.

\ct

\pov{-}

Vài ngày sau đó, tại trụ sở văn phòng làm việc mới của tập đoàn giải trí \hll.

Kuon bước vào phòng họp với một gương mặt phờ phạc mệt mỏi rõ rệt vì thiếu ngủ sau những đêm dài ròng rã ngồi kèm cặp em gái học tập. Sự tiều tụy bất thường này lập tức thu hút sự chú ý tò mò của các thành viên đang tụ họp uống trà, đặc biệt là nhóm thế hệ thứ sáu (holoX) bấy giờ mới vừa hoàn thành xong buổi phát sóng trực tiếp.

Cô nàng bác sĩ khoa học lập dị Hakui Koyori vừa nhìn thấy cậu Tổng liền hớn hở nhảy tưng tưng lao tới ghé sát mặt hỏi dồn dập, đôi mắt sáng rực rỡ đầy phấn khích: ``Oa! Kuon-senpai ơi! Em nghe mọi người đồn đại bảo là cô bé Miyako-chan (Mizuha) yêu quý của nhà anh chuẩn bị đi tham gia kỳ thi tuyển sinh bước chân vào trường cấp 3 Yên Hòa danh giá đúng không ạ!? Cái sản phẩm hồi sinh từ thí nghiệm hóa học vĩ đại của em rốt cuộc cũng đã trưởng thành đến mức đi học cấp 3 rồi sao! Em cảm thấy vô cùng tự hào và phấn khích quá đi mất thôi dã man!''

Kuon khẽ né tránh cái sự tăng động của cô nàng sói hồng, khẽ thở dài cười đáp: ``Dạ đúng thế rồi em ạ, con bé dưới danh nghĩa Mizuha đang cùng với hai đứa em gái khác của anh ôn thi sấp mặt đây này. Anh mấy ngày nay gánh tạ kèm cặp tụi nó mà muốn nhức hết cả cái đầu luôn rồi đây này.''

Nàng tổng soái La+ Darknesss đứng khoanh hai tay trước ngực trên chiếc bục gỗ của mình, khẽ hất cằm kiêu hãnh cất giọng điệu bề trên: ``Hừm! Cái kỳ thi lớp 10 của thế giới phàm trần các ngươi thì có cái gì ghê gớm đâu chứ hả cái gã quản lý kia! Sức mạnh hắc ám tối thượng của ta thừa sức quét sạch mọi con chữ trong tích tắc cơ mà!''

Nàng phó tổng soái Takane Lui khẽ cười nhẹ bước tới cốc nhẹ vào đầu La+ một cái rõ đau để răn đe, rồi ái ngại quay sang nhìn Kuon hỏi thăm: ``Em thì lại nghe giang hồ đồn thổi bảo cái kỳ thi tuyển sinh lớp 10 công lập ở nội thành Kawachi xứ anh khốc liệt và cam go kinh khủng khiếp lắm đúng không anh Kuon? Nghe nói tỉ lệ chọi còn kinh hoàng hơn cả kỳ thi đại học nữa cơ á?''

Nàng dọn dẹp Sakamata Chloe cũng hớt hải lao tới chen vào, gương mặt lộ rõ vẻ hoang mang tột độ: ``Trời đất ơi em nghe mà nổi hết cả da gà luôn rồi đây này! Có thật là chỉ cần lỡ tay làm bài sai lệch mất có 0.25 điểm lẻ thôi là đã lập tức bị đánh trượt thẳng cẳng khỏi ngôi trường mơ ước, phải chấp nhận chịu cái số phận đi học ở mấy ngôi trường dân lập xa xôi hẻo lánh không anh Kuon ơi!? Đáng sợ vcl ra luôn ấy chứ!''

Kuon khẽ gật đầu xác nhận đầy chân thực: ``Ừm, những gì các em vừa nghe hoàn toàn không có một chút xíu nào là phóng đại quá lời đâu nhé Chloe, Lui ạ. Kỳ thi vào 10 ở chỗ anh nó thực sự là một cuộc chiến sinh tử không tiếng súng đấy. Điểm chuẩn trường Yên Hòa thường niên lúc nào cũng thuộc hàng cao chót vót chạm ngưỡng trần, thế nên chỉ cần sẩy chân một phát thôi là coi như toàn bộ công sức học tập mười hai năm đèn sách bay màu sạch sẽ ngay lập tức chứ chẳng đùa đâu.''

Nàng samurai Kazama Iroha đứng bên cạnh nghe xong cũng khẽ rùng mình, hai tay siết chặt lấy thanh kiếm gỗ bảo bối, dõng dạc lên tiếng cổ vũ đầy nhiệt huyết: ``Nếu đã cam go khốc liệt đến thế thì em xin được phép gửi lời chúc may mắn, cầu chúc cho cả ba cô em gái của anh Kuon-senpai sẽ luôn giữ vững được tinh thần chiến đấu quả cảm của một samurai để chém bay mọi đề thi hóc búa, rực rỡ vượt qua cánh cổng Yên Hòa thành công nhé ạ!''

Toàn thể các cô gái thần tượng có mặt trong phòng họp bấy giờ cũng đồng loạt giơ nắm đấm lên trời, cùng nhau đồng thanh hô to gửi lời chúc mừng cổ vũ tiếp thêm sức mạnh cho gia đình cậu Tổng khiến cho Kuon đang mệt mỏi cũng cảm thấy được an ủi sướng khoái phần nào.

\ct

Và rồi, cái ngày phán quyết định mệnh cuối cùng cũng đã chính thức gõ cửa giáng xuống thành phố Kawachi.

Khung cảnh bên ngoài cánh cổng sắt lớn của ngôi trường THPT Yên Hòa sáng sớm hôm đó ngập tràn trong một bầu không khí vô cùng trang nghiêm, ngột ngạt và căng thẳng đến nghẹt thở. Hàng ngàn bậc phụ huynh đứng xếp hàng chen chúc kín mít hai bên vệ đường dưới cái nắng hè oi ả, đôi mắt ai nấy đều dõi theo bóng dáng của con em mình đang chậm rãi bước chân vào phòng thi với những gương mặt lo âu dồn nén sầu thảm.

Kuon và Kaigou đứng ở ngay sát rào chắn bảo vệ, hai tay khoanh trước ngực nghiêm túc dõi theo bóng dáng của ba cô em gái. Suzuki khẽ quay đầu lại vẫy vẫy tay chào tạm biệt hai người với một nụ cười rạng rỡ đong đầy sự quyết tâm, Mizuha (Miyako) thì vẫn luôn giữ nguyên cái điệu bộch thiệp, giữ khoảng cách vừa đủ và phóng ánh mắt cảnh giác bảo vệ xung quanh cô chủ nhỏ, còn Ryuuhou thì khom lưng nhai kẹo bơ bôm bốp tạo dáng gấu chó thách thức mọi đề thi.

Kaigou khẽ ghé sát tai nói nhỏ với Kuon bằng chất giọng sắc sảo: ``\vie{\hl{Aniki này, anh có nghĩ là cái mẻ ôn luyện cấp tốc mấy ngày qua của tụi mình sẽ thực sự giúp ích cho cái con bé bướng bỉnh kia đập tan được cái bức tường thi cử này không thế hả anh?}}''

Kuon khẽ nhấp một ngụm nước lọc, điềm nhiên đáp: ``\vie{\hl{Thì tất cả mọi sự chuẩn bị tốt nhất tụi mình đều đã dâng hiến hết sức lực lo liệu xong xuôi cho con bé rồi em ạ. Giờ thì chỉ còn biết trông chờ vào cái bản lĩnh thực chiến của con bé khi đối mặt với đề thi thực tế thôi chứ biết sao giờ. Cố lên nhé các em yêu của anh.}}''

Trải qua những ngày thi liên tiếp khốc liệt, từ các môn tự nhiên đại số Toán học hóc búa cho đến các môn văn chương dài dằng dặc, cả ba cô gái nhà Hikawa cuối cùng cũng đã hoàn thành xuất sắc bài thi cuối cùng của mình. Khi tiếng chuông báo hiệu kết thúc giờ làm bài vang dội lên xé toạc không gian trường học, toàn bộ thí sinh đồng loạt buông bút thở phào nhẹ nhõm một hơi dài xả xui, cuộc chiến tri thức cam go cuối cùng đã khép lại hoàn hảo.

Thế nhưng, chặng đường căng thẳng bế tắc thực sự thì lại mới chỉ bắt đầu từ cái giây phút ngóng chờ kết quả điểm số chính thức mà thôi.

\ct

\pov{Hikawa Kuon}

Một tuần lễ chờ đợi kết quả trôi qua trong một bầu không khí tĩnh lặng và nghẹt thở vô cùng. Ngày hôm nay chính là cái thời điểm định mệnh khi Sở giáo dục chính thức công bố bảng điểm tra cứu trực tuyến trên mạng.

Tôi cùng với Kaigou đang ngồi tề tựu trước màn hình máy tính lớn đặt ở phòng khách nhà Hikawa để chuẩn bị nhập mã số báo danh tra cứu kết quả. Đứng xếp hàng ngay ngắn ở phía sau lưng hai chúng tôi bấy giờ chính là ba cô em gái Suzuki, Ryuuhou và Mizuha đang cùng nhau nắm chặt chặt lấy bàn tay của nhau vì hồi hộp, lo sợ. Thậm chí đến cả cái con bé Ryuuhou thường ngày vốn dĩ vô cùng ngạo nghễ bướng bỉnh, bấy giờ cái gương mặt cũng đã tái mét lại vì hoảng sợ, răng cắn môi liên tục không dám hé răng nói năng câu nào sút hột.

Cô bé Naomi 10 tuổi khẽ ôm chặt lấy cánh tay tôi an ủi, cất chất giọng ngọt ngào: ``Sẽ không sao đâu Onii-chan ơi, em tin chắc chắn là các chị yêu quý của em đều sẽ đạt được kết quả tuyệt hảo nhất.''

Tôi khẽ gật đầu gật đầu nhẹ nhàng xoa đầu cô bé lấy lại sự bình tĩnh, rồi dùng tay nhanh nhẹn gõ bàn phím nhập mã tra cứu kết quả cho thí sinh đầu tiên — Hikawa Mizuha.

Giao diện trang web load xoay tròn mất vài giây, rồi bảng điểm số hiện ra rõ nét trước mắt mọi người. Kaigou đứng bên cạnh nhìn thấy con số hiển thị liền lập tức trợn tròn mắt kinh ngạc, thốt lên đầy phấn khích: ``Quả thực không ngoài dự đoán của tụi mình tí nào cả luôn nhé Aniki! Mii đạt mốc điểm số tuyệt đối tối đa 40 trên 40 điểm luôn kìa! Đúng là quái vật thủ khoa thiên tài thực thụ không ai làm lại được luôn rồi cơ!''

Miyako nghe thấy điểm số tối đa của mình nhưng gương mặt vẫn giữ nguyên nét tỉnh bơ thản nhiên hệt như không có chuyện gì, cô nàng chỉ khẽ gập người cúi đầu chàoch thiệp đáp lễ: ``Dạ, kết quả thế này là bình thường thôi ạ, em cảm ơn mọi người nhiều nha.''

Tôi khẽ mỉm cười nhẹ nhõm, rồi tiếp tục thao tác nhập mã số báo danh cho thí sinh thứ hai — Hikawa Suzuki.

Màn hình hiển thị điểm số load xong xuôi, con số tổng điểm hiển thị rõ ràng: 36/40 điểm. (Biết rằng năm tuyển sinh năm nay điểm chuẩn đầu vào lấy chuẩn của trường Anwa công bố lấy vừa vặn chạm mốc 34 điểm). Như vậy, cô em gái Suzuki của chúng ta đã xuất sắc đỗ đạt vượt qua điểm chuẩn Yên Hòa dư dả hẳn 2 điểm tròn trịa!

Suzuki nhìn thấy con số điểm đỗ đạt của mình hiển thị rực rỡ trên màn hình liền không kìm chế nổi sự xúc động nghẹn ngào, cô nàng oà khóc thút thít vì hạnh phúc tột cùng, nhảy bổ vào lòng ôm chầm lấy Kaigou: ``Hu hu chị Kaigou ơi em làm được rồi! Em đỗ đạt vào trường Yên Hòa rồi chị ơi! Công sức cần cù học tập của em cuối cùng đã được đền đáp xứng đáng rồi chị ơi!''

Kaigou cũng khẽ mỉm cười dịu dàng ôm chặt lấy cô em gái nhỏ vào lòng sực xoa đầu cưng nựng: ``Ngoan nào cưng, chị biết là em nhất định sẽ làm được. Chúc mừng em gái yêu quý của chị nhé!''

Bầu không khí trong phòng khách lúc này bỗng chốc trở nên nghẹt thở và căng thẳng tột độ khi chỉ còn sót lại duy nhất một cái tên cuối cùng chưa được tra cứu — Hikawa Ryuuhou. Toàn thân con bé bấy giờ run lên lẩy bẩy như cành cây trước dông bão, mồ hôi hột tuôn ra như tắm đầm đìa cả khuôn mặt.

Tôi nuốt nước bọt cái ực để lấy lại bình tĩnh, rồi run rẩy gõ nốt mã số báo danh của Ryuuhou nhấn Enter tra cứu.

Vòng tròn trang web xoay đều xoay đều trong một sự im lặng tĩnh mịch đến nghẹt thở hệt như đang đi tàu lượn siêu tốc của cả gia đình. Và rồi, con số tổng điểm cuối cùng đã chính thức hiển thị rõ nét trên màn hình: 34 / 40 điểm.

Một sự im lặng đứng hình bao trùm khắp căn phòng mất vài giây toàn tập. Rồi ngay sau đó, cô em gái Ryuuhou bỗng dưng trợn ngược đôi mắt lên nhìn bảng điểm, hai tay giơ cao nắm đấm lên trời nhảy dựng người lên hét lớn một tiếng vang dội chấn động cả ngôi nhà:

``\textbf{CÁI ĐM NÓ CHỨ BÀ MÀY ĐỖ RỒI NHA MẤY ĐỨA ƠI!!! BÀ MÀY ĐỖ VỪA VẶN VỪA KHÍT ĐÚNG CHUẨN ĐIỂM CHUẨN 34 ĐIỂM ĐỂ BƯỚC CHÂN VÀO TRƯỜNG YÊN HÒA RỒI NHÉ!!! CHÚNG MÀY THẤY SỨC MẠNH CỦA GAME THỦ TOXIC CHƯA HẢ CÁI LŨ ĐỒNG NÁT KIA!!! HẮC HẮC SƯỚNG KHÓAI QUÁ ĐI MẤT THÔI!!!}''

Cả gia đình nhà Hikawa bấy giờ mới hoàn toàn thở phào nhẹ nhõm một hơi dài sướng khoái trút bỏ được toàn bộ gánh nặng lo âu bấy lâu nay. Mọi người xúm lại ôm chầm lấy nhau reo hò chúc mừng nồng nhiệt cho cái sự đỗ đạt sát nút đầy ảo ma của cô em gái út bướng bỉnh.

Tôi khẽ lôi cuốn sổ ghi chép điểm số của năm xưa mình ra đối chiếu so sánh, khẽ nhếch khóe môi cười nhạt tự nhủ đầy hoài niệm: ``Chà, nghĩ lại cái năm tuyển sinh của mình ngày trước thang điểm là 70 điểm, phải thi cật lực tới tận 5 môn học gồm Toán-Văn nhân đôi hệ số, cộng thêm tiếng Anh, tiếng Nhật vàch sử bốc thăm ngẫu nhiên, thế mà mình lúc đó đỗ đạt cũng chỉ thừa có đúng 0.5 điểm lẻ thôi chứ. Còn cái mẻ ba đứa tụi em năm nay thi có 4 môn Toán-Văn-Anh-Nhật tính theo thang điểm 40 thôi, mà Mizuha thì giật điểm số tối đa 40 điểm thủ khoa, Suzuki thì dư 2 điểm, còn Ryuuhou thì vừa khít khịt đúng chuẩn 34 điểm đỗ đạt thế này thì đúng là cái kịch bản gia đình học bá Yên Hòa xuất sắc vô địch thiên hạ rồi còn đâu nữa cơ chứ.''

Ryuuhou nhí nhảnh ôm chặt lấy cổ tôi đu đưa nũng nịu làm nũng: ``Hì hì, thế tối hôm nay ông anh già có lòng tốt dắt cả nhà tụi em đi ăn một bữa tiệc lẩu sukiyaki thịt bò siêu to khổng lồ để đáp lễ chúc mừng chiến thắng cho tụi em không thế hả Onii-san yêu quý của em ơi\~{}?''

Tôi mỉm cười nhẹ nhàng xoa xoa mái tóc mềm mại của con bé gật đầu dứt khoát: ``Được rồi, được rồi, hôm nay cả nhà tụi em lập công lớn thế thì anh đây không ngại ngần gì mà xuất quỹ chiêu đãi cả nhà một bữa no say ra trò luôn nhé các em yêu của anh!''

Tiếng cười đùa vui vẻ, sảng khoái và hạnh phúc tột cùng của đại gia đình 6 anh chị em nhà Hikawa lại một lần nữa vang dội rộn rã lên khắp tòa nhà, xua tan hoàn toàn đi mọi sự mệt mỏi căng thẳng của những ngày thi cử dài đằng thẵng vừa qua. Cánh cổng ngôi trường cấp 3 Yên Hòa danh giá bấy giờ đang dang rộng vòng tay chào đón những người chủ nhân mới của nó bước vào một chuyến hành trình thanh xuân mới đầy rực rỡ và hứa hẹn phía trước.

\begin{center}
    \textbf{\Large --- HẾT ---}
\end{center}

\end{document}